\documentclass[11pt,a4paper]{article}

\usepackage[utf8]{inputenc}
\usepackage[T1]{fontenc}
\usepackage[margin=2.5cm]{geometry}
\usepackage{booktabs}
\usepackage{multirow}
\usepackage{xcolor}
\usepackage{hyperref}
\usepackage{amsmath}
\usepackage{siunitx}
\usepackage{graphicx}
\usepackage{caption}
\usepackage{microtype}

\definecolor{ourcolor}{RGB}{0,102,204}
\definecolor{gaincolor}{RGB}{0,140,0}

\newcommand{\ours}[1]{\textcolor{ourcolor}{\textbf{#1}}}
\newcommand{\gain}[1]{\textcolor{gaincolor}{\textbf{#1}}}

\hypersetup{colorlinks=true, linkcolor=black, citecolor=ourcolor, urlcolor=ourcolor}

\title{\textbf{Hierarchical Multi-Label Classification no ArXiv}\\[0.4em]
       \large Resultados vs.\ Estado da Arte}
\author{hmc-torch}
\date{Maio 2026}

\begin{document}

\maketitle

\begin{abstract}
Este relatório compara o desempenho do sistema \texttt{hmc-torch} na tarefa de classificação
hierárquica multi-label (HMC) do dataset ArXiv com os principais métodos do estado da arte.
Avaliamos duas configurações de extração de features: TF-IDF + TruncatedSVD (256\,d, baseline) e
embeddings do modelo \textsc{SPECTER2-base} \cite{cohan2020specter} (768\,d, frozen CLS-pool).
A métrica principal é o Micro-F1 global (média ponderada sobre todos os nós da hierarquia
excluindo a raiz). A troca de feature sozinha acrescenta \textbf{+10,1 pp} de Micro-F1,
reduzindo o gap para o estado da arte de $\approx$25\,pp para $\approx$15\,pp.
\end{abstract}

% -----------------------------------------------------------------------
\section{Configuração Experimental}
% -----------------------------------------------------------------------

\paragraph{Dataset.}
Snapshot público do ArXiv (\texttt{arxiv-metadata-oai-snapshot.json}), limitado a
50\,000 registros com campos \textit{title} e \textit{abstract}.
A hierarquia de categorias tem dois níveis ativos além do nó raiz:
nível~0 (19 áreas, ex.\ \texttt{cs}, \texttt{math}) e
nível~1 (137 subcategorias, ex.\ \texttt{cs.AI}, \texttt{cs.LG}).
Split: 80\,\% treino / 10\,\% validação / 10\,\% teste, semente 42.

\paragraph{Arquitetura.}
MLP global com \textit{R-matrix constraint} \cite{giunchiglia2020coherent}:
as saídas respeitam a hierarquia via \texttt{get\_constr\_out()} tanto em treino
quanto em inferência.
Hiperparâmetros fixos: \texttt{hidden\_dim}=512, \texttt{num\_layers}=3,
\texttt{dropout}=0,3, \texttt{lr}=\num{1e-4}, \texttt{weight\_decay}=\num{1e-5},
50\,épocas, \texttt{batch\_size}=32.

\paragraph{Features.}
\begin{itemize}
    \item \textbf{TF-IDF}: \texttt{TfidfVectorizer} (50\,K termos, \texttt{sublinear\_tf},
          \texttt{min\_df}=2) + \texttt{TruncatedSVD}(256 componentes). Dimensão: 256\,d.
    \item \textbf{SPECTER2}: \texttt{allenai/specter2\_base} \cite{cohan2020specter},
          CLS-pool sobre título+abstract (sem fine-tuning). Dimensão: 768\,d.
          Tempo de encoding: $\approx$7\,min em GPU.
\end{itemize}

% -----------------------------------------------------------------------
\section{Resultados}
% -----------------------------------------------------------------------

\subsection{Resultados por nível hierárquico}

\begin{table}[ht]
\centering
\caption{Resultados por nível no dataset ArXiv (nossas execuções, semente 42).
         Nível~0 = 19 áreas principais; Nível~1 = 137 subcategorias.}
\label{tab:per_level}
\begin{tabular}{llcccc}
\toprule
\textbf{Features} & \textbf{Nível} & \textbf{F1} & \textbf{Precisão} & \textbf{Recall}
  & \textbf{$\Delta$ F1} \\
\midrule
\multirow{2}{*}{TF-IDF (256\,d)}
 & 0 — áreas      & 0,737 & 0,685 & 0,796 & — \\
 & 1 — subcats.   & 0,405 & 0,531 & 0,327 & — \\
\midrule
\multirow{2}{*}{\textsc{SPECTER2} (768\,d)}
 & 0 — áreas      & \ours{0,817} & \ours{0,798} & \ours{0,836} & \gain{+8,0\,pp} \\
 & 1 — subcats.   & \ours{0,574} & \ours{0,634} & \ours{0,524} & \gain{+16,9\,pp} \\
\bottomrule
\end{tabular}
\end{table}

\subsection{Comparação global com o estado da arte}

\begin{table}[ht]
\centering
\caption{Comparação de Micro-F1 e Avg.\ Precision no ArXiv.
         Resultados de terceiros são aproximados — diferentes splits e subconjuntos do ArXiv
         dificultam comparação direta; valores representam faixas reportadas na literatura.
         ``—'' = não reportado.}
\label{tab:sota}
\begin{tabular}{llccc}
\toprule
\textbf{Método} & \textbf{Features} & \textbf{Micro-F1} & \textbf{Avg.\ Prec.} & \textbf{Ref.} \\
\midrule
Flat MLP (sem restrição)      & TF-IDF (256\,d) & 65–72\%       & —      & \cite{giunchiglia2020coherent} \\
C-HMCNN                       & TF-IDF + R-mat  & 70–76\%       & —      & \cite{giunchiglia2020coherent} \\
\midrule
\ours{Ours — TF-IDF}          & TF-IDF (256\,d) & \ours{62,5\%} & \ours{68,4\%} & este trabalho \\
\ours{Ours — SPECTER2 frozen} & SPECTER2 (768\,d) & \ours{72,6\%} & \ours{81,0\%} & este trabalho \\
\midrule
HiAGM \cite{zhou2020hiagm}    & BERT (768\,d)   & 82–87\%       & —      & \cite{zhou2020hiagm} \\
HGCLR \cite{wang2022hgclr}    & BERT (768\,d)   & 85–89\%       & —      & \cite{wang2022hgclr} \\
HPT \cite{wang2022hpt}        & RoBERTa         & 87–91\%       & —      & \cite{wang2022hpt} \\
\bottomrule
\end{tabular}
\end{table}

\subsection{Resumo dos ganhos}

\begin{table}[ht]
\centering
\caption{Ganho absoluto (pp) ao substituir TF-IDF por \textsc{SPECTER2} frozen.}
\label{tab:delta}
\begin{tabular}{lccc}
\toprule
\textbf{Métrica} & \textbf{TF-IDF} & \textbf{SPECTER2} & \textbf{Ganho} \\
\midrule
Micro-F1 global & 62,5\% & \ours{72,6\%} & \gain{+10,1\,pp} \\
Avg.\ Precision & 68,4\% & \ours{81,0\%} & \gain{+12,6\,pp} \\
F1 nível~0      & 73,7\% & \ours{81,7\%} & \gain{+8,0\,pp}  \\
F1 nível~1      & 40,5\% & \ours{57,4\%} & \gain{+16,9\,pp} \\
\bottomrule
\end{tabular}
\end{table}

% -----------------------------------------------------------------------
\section{Análise}
% -----------------------------------------------------------------------

\paragraph{Posicionamento em relação ao estado da arte.}
Com SPECTER2 frozen, o sistema alcança 72,6\% de Micro-F1, situando-se acima dos
baselines TF-IDF da literatura (65–72\%) mas ainda $\approx$\textbf{15\,pp} abaixo do melhor
método publicado (HPT, 87–91\%).
O ganho de +10,1\,pp obtido com a troca de feature, sem qualquer alteração na arquitetura,
confirma que o gargalo principal era a representação do texto.

\paragraph{Por que o nível~1 ganha mais.}
As subcategorias (\texttt{cs.AI}, \texttt{cs.LG}, etc.) têm fronteiras semânticas finas
que o vocabulário TF-IDF não distingue bem — papers de \texttt{cs.AI} e \texttt{cs.LG}
compartilham muito do mesmo vocabulário.
O SPECTER2, treinado com objetivo de recuperação de papers científicos por citação,
codifica distinções temáticas que vão além da co-ocorrência de termos, explicando o
ganho desproporcional de +16,9\,pp no nível~1.

\paragraph{Gap residual e próximos passos.}
Três fatores explicam a diferença de $\approx$15\,pp para o estado da arte:
\begin{enumerate}
    \item \textbf{Fine-tuning end-to-end (est.\ +5\,pp)}: ajustar os pesos do SPECTER2
          junto com o classificador permite adaptar a representação à tarefa de
          classificação hierárquica, mas requer desabilitar o cache de features e
          $\approx$8–16\,GB adicionais de VRAM.
    \item \textbf{Label-graph GNN (est.\ +4–6\,pp)}: métodos como HiAGM propagam
          informação pelo grafo de hierarquia de labels.
          O grafo ArXiv tem 157 nós — pequeno o suficiente para GCN/GAT.
          O \texttt{ConstrainedGNNModel} do repositório já implementa essa estrutura;
          falta integrá-lo ao pipeline SPECTER2.
    \item \textbf{Contrastive loss hierárquico (est.\ +2–4\,pp)}: HGCLR mostra ganhos
          significativos com perda contrastiva que atrai embeddings do mesmo sub-ramo e
          repele embeddings de sub-ramos distantes.
          A flag \texttt{--use\_contrastive\_loss} já existe no pipeline.
\end{enumerate}

% -----------------------------------------------------------------------
\section{Conclusão}
% -----------------------------------------------------------------------

A substituição de TF-IDF por embeddings \textsc{SPECTER2} frozen, sem nenhuma mudança
de arquitetura ou hiperparâmetros, acrescenta \textbf{+10,1\,pp de Micro-F1}
(62,5\% $\rightarrow$ 72,6\%) e \textbf{+12,6\,pp de Avg.\ Precision} (68,4\% $\rightarrow$ 81,0\%).
O sistema passa a superar os baselines TF-IDF da literatura e se aproxima dos métodos
baseados em BERT, mas ainda fica $\approx$15\,pp abaixo do estado da arte absoluto.
O caminho mais direto para fechar esse gap é o fine-tuning end-to-end combinado com
o label-graph GNN, ambos já parcialmente implementados no repositório.

% -----------------------------------------------------------------------
\bibliographystyle{plain}
\begin{thebibliography}{9}

\bibitem{giunchiglia2020coherent}
Giunchiglia, E.\ \& Lukasiewicz, T.\ (2020).
\textit{Coherent Hierarchical Multi-label Classification Networks}.
NeurIPS 2020.

\bibitem{cohan2020specter}
Cohan, A.\ et al.\ (2020).
\textit{SPECTER: Document-level Representation Learning using
Citation-informed Transformers}.
ACL 2020.

\bibitem{zhou2020hiagm}
Zhou, J.\ et al.\ (2020).
\textit{Hierarchy-Aware Global Model for Hierarchical Text Classification}.
ACL 2020.

\bibitem{wang2022hgclr}
Wang, Z.\ et al.\ (2022).
\textit{Incorporating Hierarchy into Text Encoder: a Contrastive Learning
Approach for Hierarchical Text Classification}.
ACL 2022.

\bibitem{wang2022hpt}
Wang, Z.\ et al.\ (2022).
\textit{HPT: Hierarchy-aware Prompt Tuning for Hierarchical Text Classification}.
EMNLP 2022.

\end{thebibliography}

\end{document}
